\documentclass[10pt]{article}
\usepackage[letterpaper,margin=0.8in]{geometry}
\usepackage{amsmath,amssymb,booktabs,array}
\usepackage[T1]{fontenc}
\usepackage{lmodern}
\usepackage[hidelinks]{hyperref}
\usepackage{xurl}
\setlength{\parindent}{0pt}
\setlength{\parskip}{4pt}
\makeatletter
\renewcommand{\maketitle}{\begin{center}{\Large\bfseries\@title\par}\vspace{5pt}\@author\qquad\@date\end{center}}
\makeatother
\title{Scope Dependence of Strict Splitting Counts\\
\large An Exact Check of 69 Recorded Seed Positions}
\author{Anonymous}
\date{30 October 2026}
\begin{document}
\maketitle
\begin{abstract}
We test the previously reported lower strict-splitting-count criterion on a
broader, previously certified census of 69 prime-order seed positions, retaining
all 82 recorded root branches. No cutoff from 100 through 5000 separates the
ten published positions from all 59 unpublished positions. At the original
cutoff 691, 22 unpublished positions tie or fall below their field's published
maximum. They belong to categories excluded from the earlier 21-choice cohort
by additional root and order-width screens. A separate PARI implementation
reproves 76 distinct field/order primes, verifies the point certificates and
reconstructs all six output files; three deliberate corruptions are refused.
The earlier 21-choice result reproduces exactly. This scope test limits the
splitting criterion alone; it neither refutes a combined-screen hypothesis nor
establishes historical selection, random-curve unusualness or an ECDLP weakness.
\end{abstract}

\section{Fixed criterion and population}
For a prime-field curve of certified full prime order $n$, define
\[
 t=p+1-n,\quad \Delta=t^2-4p,\quad
 S_L=\#\{\text{odd primes }\ell\le L:(\Delta/\ell)=+1\}.
\]
Let $R_L$ count the separate zero-character cases. A strict split prime has
two distinct Frobenius-stable lines in $E[\ell]$, corresponding to cyclic
isogeny kernels defined over the base field. Their individual points need
not be rational. Ramified cases can also admit rational isogenies; $S_L$
therefore does not count every rational small-degree isogeny~\cite{ss,pari}.

The degree-691 observation predates this analysis~\cite{asanso}. The earlier
21-choice replay compares ten publications with eleven named unpublished
candidates. Their lower-count separation survives only at integer cutoffs
691--718 and 769--772. The inherited November-1997-draft acceptance analysis
records a broader prime-order cohort: 69 positions, of which ten are published
and 59 unpublished~\cite{acceptance}. Its additional 48 positions comprise
17 width-only and 31 other-root-only cases. This is a lower-bound census of
numerically eligible recorded positions, not evidence of original visits,
rejections or an exhaustive near-prime search.

The fixed comparison is the same: within each field, all unpublished counts
must exceed the largest published count. At each integer $100\le L\le5000$,
the margin is $m_L=\min_{\rm unpub}S_L-\max_{\rm pub}S_L$; strict success
requires $m_L>0$. The broader census has five paired fields. P-521 remains
a singleton and receives no artificial margin. All labels were visible;
this is retrospective reanalysis, with no fitted replacement cutoff.

\section{Broader result}
\begin{center}\small\begin{tabular}{rrrrrrrr}\toprule
Bits & Pub. & Unpub. & Pub. $S_{691}$ & Unpub. $S_{691}$ & Margin & Pos. cutoffs & Violators\\\midrule
192 & 3 & 23 & 57--60 & 44--71 & -16 & 0/4901 & 14\\
224 & 1 & 11 & 50--50 & 50--66 & 0 & 1801/4901 & 1\\
239 & 3 & 14 & 42--57 & 54--71 & -3 & 0/4901 & 5\\
256 & 1 & 5 & 57--57 & 56--67 & -1 & 0/4901 & 2\\
384 & 1 & 6 & 54--54 & 56--60 & 2 & 635/4901 & 0\\
521 & 1 & 0 & 52--52 & --- & --- & --- & ---\\
\bottomrule\end{tabular}
\end{center}
The ranges refer to $S_{691}$. Positive-cutoff counts cover all 4901 integers.
``Violators'' are unpublished positions with $S_{691}$ at or below the
published maximum in the same field; equality fails the strict rule.
No broader joint strict or inclusive ($S_L+R_L$) separator occurs anywhere
in the requested range. P-192, P-239 and P-256 each fail individually at
every cutoff. This does not change the correctly reproduced narrow result.
\clearpage

\section{Counterexamples and the role of additional screens}
\begin{center}\small\begin{tabular}{lrr}\toprule
Recorded category & Positions & Violators at691\\\midrule
Other-root-only & 31 & 14\\
Order-width-only & 17 & 8\\
Named omissions & 11 & 0\\
\bottomrule\end{tabular}
\end{center}
All 22 counterexamples lie among the extra 48 positions. The eleven named
omissions remain above their field's published maximum at 691. Accordingly,
the broader failure limits \emph{splitting alone}; a conjunction that first
applies the root/width screens and then the splitting test still fits these
recorded labels. Its historical use and motivation remain unverified. Treating
the extra positions as original rejections would also require unavailable
visitation evidence. The table gives every strict-691 counterexample.

\begin{center}\small\begin{tabular}{lcrrr}\toprule
Seed position & Category & $S_{691}$ & Pub. maximum & Difference\\\midrule
p192-338 & O & 58 & 60 & -2\\
p192-774 & O & 60 & 60 & 0\\
p192-1006 & W & 56 & 60 & -4\\
p192-1256 & O & 58 & 60 & -2\\
p192-1477 & O & 59 & 60 & -1\\
p192-1914 & O & 58 & 60 & -2\\
p192-2678 & W & 55 & 60 & -5\\
p192-2771 & W & 44 & 60 & -16\\
p192-4254 & O & 57 & 60 & -3\\
p192-5242 & W & 58 & 60 & -2\\
p192-5674 & W & 52 & 60 & -8\\
p192-5701 & W & 54 & 60 & -6\\
p192-5964 & W & 56 & 60 & -4\\
p192-5976 & O & 51 & 60 & -9\\
p224-2328 & W & 50 & 50 & 0\\
p239-155 & O & 56 & 57 & -1\\
p239-765 & O & 56 & 57 & -1\\
p239-1670 & O & 55 & 57 & -2\\
p239-2228 & O & 56 & 57 & -1\\
p239-2311 & O & 54 & 57 & -3\\
p256-276 & O & 56 & 57 & -1\\
p256-1383 & O & 56 & 57 & -1\\
\bottomrule\end{tabular}
\end{center}
Category W denotes order-width-only; O denotes other-root-only. The difference
is the candidate's score minus the largest published score in its field.

\section{Branches, the base point and coefficient}
All 82 recorded branches remain explicit. Within each seed position, their
discriminants are verified equal before selecting a representative for the
69-position comparison: canonical when present, otherwise the sole other
branch. Isomorphic P-224 coefficient signs have equal order. A quadratic
twist over the other fields has trace $-t$, and $(-t)^2-4p=\Delta$.
Thus this score cannot distinguish those twist branches.

For $a=-3$, every on-curve point satisfies
$b\equiv y^2-x^3+3x\pmod p$. Hence $x=0$ gives $y^2=b$, while $x=1$
gives $y^2=b-2$. These identities explain simple relations for assigned
smallest-$x$ points. The earlier finite G--b study independently checked
100548 trait cells, found no separator among 19152 single-envelope tests,
and no match in 880 fixed public-scalar checks~\cite{gb}. It does not
exhaust all mathematical functions, and the historical unpublished generators
remain unknown. The present score uses only $p,n$ and adds no point attack.
\clearpage

\section{Verification, provenance and interpretation}
Source, inputs, commands and resource bounds were committed before actual
cohort extraction at \path{1fe0d49974b2a3083772d003b69f4e75fe48aa4b}.
The GMP producer checks nonsingularity, point equations, Hasse membership,
uniqueness and $[n]P=O$, then enumerates quadratic residues. The separate PARI
verifier uses proved \texttt{isprime}, \texttt{ellmul} and \texttt{kronecker},
and pairwise comparison differences. It reproves all 76 distinct field/order
primes and reconstructs six files byte for byte: branch endpoints, characters,
all seed/cutoff profiles, joint outcomes/margins, counterexamples and summary.
The implementations share the GMP integer family and TSV I/O, which handles
no curve or comparison arithmetic. No new production point count was needed.

All ten phases finish as expected: two arithmetic controls, producer, verifier
and three corruption preparations exit zero; the three deliberately corrupted
character/profile/cutoff copies are refused with native exit one. There are
54776 odd-prime branch characters, 338169 seed/cutoff profiles and 4901 joint
cutoff outcomes. Every original branch, source file, phase and corrupted copy
is retained in the local raw archive. No S3 upload or remote custody receipt
is claimed. Existing-results analysis uses the repository's documented local
development exception; no random-experiment gate is waived.

The finding is a restriction on the empirical criterion's scope. It does
not establish the original decision program, either weak-publication or
weak-omission hypothesis, or a designer's motive. Fewer strict split primes
has no demonstrated ECDLP-hardness ordering on these inputs. Isogeny-based
random-reducibility results have additional same-field/order/endomorphism
assumptions and do not turn this count into an attack~\cite{jmv}.
The requested 10000 accepted random prime-order curves per field and published
percentiles remain unrun. Already examined large fields are not unseen
holdouts. A finite test family cannot measure every mathematical difference.

\begin{thebibliography}{9}\small
\bibitem{asanso} asanso. Public strict splitting-count observation, 2026.
Pinned provenance in the earlier strict-count study,
\path{research/nist_skip_criterion_20261025/strict_elkies_20261029}.
\bibitem{acceptance} Prior lab analysis. November-1997 X9.62 acceptance-space
replay, 24 October 2026. Source revision
\path{351c9c021c7be044ade54dc7e352b33d3a772244},
\path{research/nist_seed_1997_acceptance_20261024/analysis.json}.
\bibitem{ss} I. Shparlinski and A. Sutherland. On the distribution of Atkin
and Elkies primes. \url{https://arxiv.org/abs/1112.3390}, version 2.
\bibitem{jmv} D. Jao, S. Miller and R. Venkatesan. Expander graphs based on
GRH with an application to elliptic curve cryptography.
\url{https://arxiv.org/abs/0811.0647}, version 2.
\bibitem{pari} PARI Group. PARI/GP documentation: \texttt{isprime},
\texttt{kronecker}, \texttt{ellinit}, \texttt{ellmul}.
\url{https://pari.math.u-bordeaux.fr/dochtml/html/}.
\bibitem{gb} Prior lab study. Generator--coefficient comparison: all 21 curves,
29 October 2026. \path{research/nist_skip_criterion_20261025/generator_b_execution_20261029/REPORT.md}.
\end{thebibliography}
\clearpage
\appendix
\section{All 69 selected seed-position endpoints}
P: published; N: named unpublished omission; W: order-width-only;
O: other-root-only. These are source classification categories, not proven
historical program dispositions. $R$ is retained separately. Representative
choice and both branch endpoints are explicit in the raw 82-branch table.
Full profiles include all cutoffs, not only these endpoints.

\begin{center}\small\begin{minipage}[t]{0.48\linewidth}\centering
\begin{tabular}{lcrrr}\toprule
Seed position & Cat. & $S_{691}$ & $R_{691}$ & $S_{5000}$\\\midrule
p192-139 & P & 60 & 3 & 335\\
p192-338 & O & 58 & 2 & 323\\
p192-774 & O & 60 & 4 & 354\\
p192-1006 & W & 56 & 2 & 307\\
p192-1256 & O & 58 & 1 & 329\\
p192-1477 & O & 59 & 1 & 347\\
p192-1914 & O & 58 & 1 & 332\\
p192-2034 & O & 64 & 0 & 361\\
p192-2198 & N & 61 & 0 & 339\\
p192-2524 & O & 64 & 1 & 331\\
p192-2678 & W & 55 & 4 & 325\\
p192-2771 & W & 44 & 2 & 314\\
p192-2832 & O & 64 & 1 & 333\\
p192-3087 & O & 71 & 3 & 351\\
p192-3506 & O & 66 & 2 & 347\\
p192-4254 & O & 57 & 4 & 323\\
p192-4455 & O & 63 & 1 & 342\\
p192-4492 & O & 64 & 0 & 325\\
p192-5228 & P & 57 & 1 & 317\\
p192-5242 & W & 58 & 5 & 334\\
p192-5405 & N & 62 & 0 & 313\\
p192-5674 & W & 52 & 2 & 311\\
p192-5701 & W & 54 & 3 & 327\\
p192-5964 & W & 56 & 3 & 341\\
p192-5976 & O & 51 & 1 & 315\\
p192-6013 & P & 58 & 1 & 303\\
p224-77 & W & 55 & 1 & 334\\
p224-777 & N & 53 & 4 & 318\\
p224-1485 & N & 56 & 3 & 342\\
p224-2232 & W & 64 & 2 & 319\\
p224-2328 & W & 50 & 1 & 314\\
p224-2938 & W & 54 & 0 & 329\\
p224-2954 & W & 59 & 3 & 319\\
p224-3191 & W & 61 & 2 & 339\\
p224-3199 & W & 65 & 4 & 340\\
\bottomrule\end{tabular}\end{minipage}
\hfill
\begin{minipage}[t]{0.48\linewidth}\centering
\begin{tabular}{lcrrr}\toprule
Seed position & Cat. & $S_{691}$ & $R_{691}$ & $S_{5000}$\\\midrule
p224-5736 & N & 59 & 2 & 325\\
p224-5861 & W & 66 & 1 & 337\\
p224-7026 & P & 50 & 3 & 299\\
p239-21 & O & 68 & 2 & 326\\
p239-36 & N & 71 & 2 & 347\\
p239-155 & O & 56 & 1 & 309\\
p239-765 & O & 56 & 1 & 328\\
p239-1280 & O & 60 & 2 & 333\\
p239-1562 & N & 67 & 3 & 337\\
p239-1670 & O & 55 & 1 & 321\\
p239-1769 & O & 68 & 1 & 325\\
p239-2185 & P & 42 & 1 & 315\\
p239-2228 & O & 56 & 4 & 346\\
p239-2285 & O & 65 & 3 & 350\\
p239-2311 & O & 54 & 1 & 327\\
p239-2320 & P & 57 & 2 & 320\\
p239-2491 & N & 63 & 2 & 352\\
p239-2682 & N & 61 & 4 & 322\\
p239-2863 & N & 66 & 1 & 329\\
p239-3048 & P & 57 & 1 & 336\\
p256-81 & O & 62 & 3 & 320\\
p256-276 & O & 56 & 5 & 318\\
p256-876 & N & 58 & 2 & 318\\
p256-1049 & O & 67 & 3 & 336\\
p256-1383 & O & 56 & 2 & 307\\
p256-1842 & P & 57 & 2 & 328\\
p384-592 & O & 58 & 1 & 320\\
p384-1432 & O & 60 & 2 & 326\\
p384-1462 & O & 56 & 1 & 332\\
p384-1497 & W & 60 & 4 & 319\\
p384-1626 & W & 60 & 1 & 308\\
p384-1796 & O & 58 & 1 & 335\\
p384-2688 & P & 54 & 0 & 324\\
p521-1136 & P & 52 & 1 & 339\\
\bottomrule\end{tabular}\end{minipage}
\end{center}
\end{document}
